% !TeX root = ../main.tex
\section{Ахах}

Путь (цепь) — незамкнутый маршрут, в котором все рёбра различны;

Цикл — замкнутый маршрут, в котором все рёбра различны;

Простой путь — путь, в котором все вершины различны;

Простой цикл — цикл, в котором все вершины различны, кроме совпадающих начала и конца.

\Def \textit{$d$-регулярный граф} - это граф, в котором каждая вершина имеет степень $d$.

\thm{Формула Кэли}{
  Количество различных деревьев на $n$ вершинах равно $n^{n-2}$.
} {
  Построим по дереву следующую последовательность: на каждой итерации будем находить лист с самым маленьким номером, запишем его родителя, а лист удалим. Получившаяся последовательность длины $n-2$, называется кодом Прюфера.
  
  Далее нужно доказать, что по кодам Прюфера однозначно восстанавливается дерево и получить биекцию между кодами  и деревьями.
}

\utv{}{
  Если $F(n, r)$ - число лесов c $r$ деревьями на $n$ вершинах, то $F(n, r) = r n^{n-1-r}$.
  Важно фиксировать корни деревьев.
}{
  Теперь обобщим код Плюфера. Пока некорневых вершин больше одной, удаляем минимальный \textbf{некорневой} лист и записываем его родителя, затем удаляя этот лист. Получаем код длины $n-r-1$, а потом записываем, к какому корню привязан последний лист.
}

Положим $u_n$ — количество унициклических графов на $n$ вершинах.

\lem{}{
  Пусть $r$ - количество вершин в цикле унициклического графа. Тогда
  \[
  u_n = \sum_{r=3}^{n} \binom{n}{r} \frac{(r-1)!}{2} \cdot F(n, r) = \sum_{r=3}^{n} \binom{n}{r} \frac{(r-1)!}{2} \cdot r \cdot n^{n-1-r}
  \]
}{
  Достаточно заметить что унициклический граф это цикл к вершинам которого прикреплены деревья.

  Тогда: Выбираем $r$ вершин цикла $\binom{n}{r}$ способами. Строим на них цикл $\frac{(r-1)!}{2}$ способами \\ ($r$ порядков вершин, r циклических сдвигов, 2 направления обхода). К каждой вершине цикла прикрепляем дерево, всего $F(n, r)$ способов.
}

\utv{}{
  $u_n \sim \sqrt{\frac{\pi}{8}} \cdot n^{n-\frac{1}{2}}$
}{
  \begin{align*}
  \binom{n}{r} &= \frac{n(n-1)\cdots(n-r+1)}{r!} =\frac{n^r}{r!} \prod_{k=0}^{r-1} \left(1 - \frac{k}{n}\right) =\\
  &= \frac{n^r}{r!} e^{\ln{(1-\frac{1}{n})}+\ln{(1-\frac{2}{n})}+\ldots+\ln{(1-\frac{r-1}{n})}} = (\star) = \frac{n^r}{r!} e^{-\frac{1+2+\ldots+(r-1)}{n}+O(\frac{r^3}{n^2})} = \frac{n^r}{r!} e^{-\frac{r(r-1)}{2n}+O(\frac{r^3}{n^2})}
  \end{align*}
  С другой стороны
  \begin{align*}
  (\star) = \frac{n^r}{r!} \cdot e^{-\frac{1}{n}-\frac{2}{n}-\ldots-\frac{r-1}{n} + O\left(\frac{1}{n^2}+\frac{4}{n^2}+\ldots+\frac{(r-1)^2}{n^2}\right)} = \frac{n^r}{r!} \cdot e^{-\frac{r(r-1)}{2n}+O(\frac{r^3}{n^2})}
  \end{align*}

  Подставим в лемму:
  \[u_n =\sum_{r=3}^{n}\left(C_n^r\frac{(r-1)!}{2}r\cdot n^{n-1-r}\right)\leq\frac{1}{2}n^{n-1}\sum_{r=3}^{n}e^{-\frac{r(r-1)}{2n}}
  \]

  \[ u_n = \sum_{r=3}^{n} \left( C_n^r \frac{(r-1)!}{2} r\cdot n^{n-1-r} \right) = \frac{1}{2} n^{n-1} \sum_{r=3}^{n} e^{-\frac{r(r-1)}{2n}} + O\left(\frac{r^3}{n^3}\right)
  \]

  \[ u_n = \sum_{r=3}^{n} \left( C_n^r \frac{(r-1)!}{2} r\cdot n^{n-1-r} \right) = \frac{1}{2} n^{n-1} \left( \underbrace{\sum_{r=3}^{[n^{0.6}]} \frac{C_n^r}{r!n^{-r}}}_{S_1} + \underbrace{\sum_{r=[n^{0.6}]+1}^{n} \frac{C_n^r}{r!n^{-r}}}_{S_2} \right)
  \]

  \[ S_2 \leq \sum_{r=[n^{0.6}]+1}^{n} e^{-\frac{r(r-1)}{2n}} \leq n\cdot e^{-\frac{n^{0.2}}{2}(1+o(1))}
  \]

  Помним, что $r\geq n^{0.6}$.
  
  \[ e^{-\frac{r(r-1)}{2n}} = e^{-\frac{r^2}{2n}(1+o(1))} \leq e^{-\frac{n^{1.2}}{2n}(1+o(1))} = e^{-\frac{n^{0.2}}{2}(1+o(1))}
  \]

  Докажем, что $\sum_{r=3}^{[n^{0.6}]} e^{-\frac{r^2}{2n}} \sim \sum_{r=0}^{\infty} e^{-\frac{r^2}{2n}}.$

  \[
  \sum_{r=0}^{2} e^{-\frac{r^2}{2n}} = O(1)
  \]

  \[
  \sum_{r=[n^{0.6}]+1}^{\infty} e^{-\frac{r^2}{2n}} = \underbrace{\sum_{r=[n^{0.6}]+1}^{n^2} e^{-\frac{r^2}{2n}}}_{(1)} + \underbrace{\sum_{r=n^2+1}^{\infty} e^{-\frac{r^2}{2n}}}_{(2)}
  \]

  Для (1):

  \[
  r\geq n^{0.6} \Rightarrow e^{-\frac{r^2}{2n}} \leq e^{-\frac{n^{2}}{2n}} = e^{-\frac n2} \Rightarrow (1) \leq n^2e^{-\frac n2}
  \]

  Для (2):

  \[
  \frac{ e^{-\frac{(r+1)^2}{2n}} }{ e^{-\frac{r^2}{2n}} } = e^{-\frac{(r+1)^2+r^2}{2n}} = e^{-\frac{2r+1}{2n}} = e^{-\frac rn-\frac{1}{2n}} < e^{-\frac rn} < e^{-n}
  \]

  Тогда (2) $\leq e^{-\frac{(n^2+1)^2}{2n}}$ (ограничили сверху геометрической прогрессией).

  Тогда:

  \begin{align*}
  S_1 &= \sum_{r=3}^{[n^{0.6}]} e^{-\frac{r(r-1)}{2n}} + O\left(\frac{r^3}{n^3}\right) \sim \sum_{r=3}^{[n^{0.6}]} e^{-\frac{r(r-1)}{2n}} \sim \sum_{r=3}^{[n^{0.6}]} e^{-\frac{r^2}{2n}} \sim \sum_{r=0}^{\infty} e^{-\frac{r^2}{2n}} \sim_{\text{сх.}} \int_0^\infty e^{-\frac{r^2}{2n}}\,dr =\\
  &= \int_0^\infty e^{-\frac{x^2}{2}} \sqrt n\,dx = \sqrt n \int_0^{+\infty} e^{-\frac{x^2}{2}}\,dx = \sqrt n \frac{ \int_{-\infty}^{+\infty} e^{-\frac{x^2}{2}}\,dx }{2} = \sqrt n \frac{\sqrt{2\pi}}{2}.
  \end{align*}
}

\Def {\textit{Клика} - это полный подграф. $\omega(G)$ - размер наибольшей клики.}

\Def {\textit{Независимое множество} - это антиклика. $\alpha(G)$ - размер наибольшего независимого множества.}